\documentclass[12pt]{article}
\usepackage[utf8]{inputenc}
\usepackage[T1]{fontenc}
\usepackage{hyperref}
\title{Sinyar Language Resources for K-12 Curriculum Development}
\author{Compiled by Claude Sonnet 5.5 (claude-sonnet-5-5) for Andr\'as Kornai}
\date{}
\begin{document}
\maketitle

\section{Introduction}
Sinyar (also Shemya or Shamya; ISO 639-3: sys; Glottocode siny1243) is a Central Sudanic language spoken around Goz Be\"ida in the Ouadda\"i region of Chad, and formerly in Darfur, Sudan. Wikipedia gives 33,000 native speakers (2023); Joshua Project gives 37,000 to 38,000 for the ethnic group. Its classification within Central Sudanic is debated. Search found academic linguistic descriptions but no educational materials.

\section{Primary Educational Resources}
No primers, textbooks or school programmes for Sinyar were found in search results. The language is not known to be taught in schools.

\section{Linguistic and Reference Materials}
\begin{itemize}
\item Wikipedia, ``Sinyar language'': overview of the language and its classification. \url{https://en.wikipedia.org/wiki/Sinyar_language}
\item Study of case marking and proper names in Sinyar (Chad/Sudan), in \emph{Language Typology and Universals} (2019). \url{https://www.degruyterbrill.com/document/doi/10.1515/stuf-2019-0019/html}
\item HAL record, ``Proper names and case markers in Sinyar (Chad/Sudan)''. \url{https://hal-supelec.archives-ouvertes.fr/LLACAN/halshs-03096955v1}
\item Chapter draft on Central Sudanic languages, which treats Sinyar-Formona as sparsely documented. \url{https://halshs.archives-ouvertes.fr/halshs-03097024v1/file/CSD_OUP%20TEXT%20%28Central%20Sudanic%29.pdf}
\item Wikipedia, ``Central Sudanic languages''. \url{https://en.wikipedia.org/wiki/Central_Sudanic_languages}
\item Numerals page for Sinyar, Ouadda\"i region. \url{https://lingweb.eva.mpg.de/channumerals/Sinyar.htm}
\end{itemize}

\section{Religious and Community Resources}
\begin{itemize}
\item Joshua Project profile of the Sinyar in Chad. \url{https://joshuaproject.net/people_groups/14941}
\end{itemize}

\section{Digital and Audio Resources}
No audio, video or corpus resources for Sinyar were found in search results.

\section{Curriculum Development Context}
Sinyar has little documentation beyond a few academic papers, and no known literacy or school materials. Speakers also use Chadian Arabic and Fur, and French and Arabic are the main languages of schooling in Chad. A K-12 curriculum would depend on new documentation, orthography development and community involvement.

\end{document}